\section{Related Work}
\label{sec:related}

\paragraph{Retrieval-augmented generation.}
RAG was introduced as a trainable combination of a dense
retriever~\cite{karpukhin2020dpr} and a sequence-to-sequence
generator~\cite{lewis2020rag}, alongside a family of retrieval-augmented
language models that condition on retrieved text during pre-training or
decoding --- REALM~\cite{guu2020realm}, FiD~\cite{izacard2021fid},
RETRO~\cite{borgeaud2022retro}, $k$NN-LM~\cite{khandelwal2020knnlm}.
Contemporary practice replaces the trained combination with a frozen
instruction-following LLM reading retrieved passages in its
prompt~\cite{ram2023incontext,shi2024replug}, the form studied here. Surveys chart the field's rapid growth
and its progression from naive through advanced to modular
pipelines~\cite{gao2023ragsurvey} and, most recently, to agentic
systems~\cite{singh2025agenticrag}.

\paragraph{Retrieval-side refinements.}
Classical lexical scoring~\cite{robertson2009bm25} survives inside modern
stacks as the complement to dense retrieval, with Reciprocal Rank
Fusion~\cite{cormack2009rrf} the standard way to merge the two rankings.
Cross-encoder reranking~\cite{nogueira2019passage} re-scores a dense
shortlist with a jointly-encoded relevance model. These methods keep one
generator call per question and modify only which passages reach it.

\paragraph{Query-side refinements.}
HyDE~\cite{gao2022hyde} searches with a model-drafted hypothetical answer
passage instead of the question, bridging the vocabulary gap between
questions and encyclopedic text. IRCoT~\cite{trivedi2023ircot} interleaves
one sentence of chain-of-thought reasoning~\cite{wei2022cot} with one
retrieval per step, so that entities surfaced by reasoning become
retrievable queries --- the canonical mechanism for multi-hop questions,
related in spirit to acting-and-reasoning agents~\cite{yao2023react}.
Neighbouring iterative schemes decompose the question into explicit
sub-questions (Self-Ask~\cite{press2023selfask}), retrieve whenever the
model's next sentence is uncertain (FLARE~\cite{jiang2023flare}), or
alternate full generation and retrieval rounds
(Iter-RetGen~\cite{shao2023iterretgen}).

\paragraph{Self-assessing and adaptive systems.}
Self-RAG~\cite{asai2024selfrag} trains a model to emit reflection tokens
deciding when to retrieve and whether generations are supported. CRAG
\cite{yan2024crag} grades retrieved passages with a trained evaluator and
triggers corrective web search when the retrieval is judged poor.
Adaptive-RAG~\cite{jeong2024adaptiverag} trains a classifier that routes
each question to no-retrieval, single-step, or multi-step pipelines. All
three decide how much work a question needs; the cascade proposed here
differs in \emph{when} the decision is made --- after the reader has seen
the evidence (a posterior signal) rather than from the question alone (a
prior one) --- and in requiring no trained component at all.

\paragraph{Structure-augmented retrieval.}
GraphRAG~\cite{edge2024graphrag}, RAPTOR~\cite{sarthi2024raptor} and
HippoRAG~2~\cite{gutierrez2025hipporag2} preprocess the corpus into graphs,
trees, or memory structures. They occupy a different cost class --- their
indexing step itself involves large-scale LLM calls over the corpus --- and
systematic evaluation finds their advantage concentrated in query-focused
summarisation rather than factoid QA~\cite{han2025ragvsgraphrag}; we
therefore exclude them from a comparison whose corpus is 21M passages, and
note the exclusion as a limitation (Section~\ref{sec:limitations}).

\paragraph{Benchmarks and evaluation.}
FlashRAG~\cite{jin2024flashrag} standardises datasets, corpora, and
reimplementations of a dozen RAG methods, and is the source of our corpus
and question sets; MultiHop-RAG~\cite{tang2024multihoprag} benchmarks
multi-hop retrieval specifically, BEIR~\cite{thakur2021beir} zero-shot
retrieval, and RGB~\cite{chen2024benchmarking} the generation step's
robustness to noisy passages; RAGAS~\cite{es2024ragas} offers
reference-free evaluation, and best-practice studies tune pipeline
components in isolation~\cite{wang2024bestpractices}. Our work differs
from toolkit
benchmarking in three ways: paired significance testing with multiplicity
control instead of point scores; measurement of the corpus-imposed recall
ceiling and of abstention as a behaviour distinct from error; and a
per-question failure diagnosis that connects each architecture's gain to
the failure mode it repairs.

\paragraph{Cascades.}
Deferring expensive computation until a cheap stage proves insufficient has
a long history --- boosted detector cascades in
vision~\cite{viola2001cascade} and cascade ranking in
retrieval~\cite{wang2011cascade}. The architecture proposed here applies
the idea at the level of whole RAG strategies, with the reader's own
refusal as the deferral signal; to our knowledge the losslessness of that
trigger under span-level Exact Match has not previously been stated or
exploited.
